\section{Plugin Protocol}
\label{sec:plugin-protocol}

\subsection{Why a plugin protocol in v1}

PQCrypt's core spec does \emph{not} define a key-on-disk format. Specifying one in v1 would require choosing a passphrase KDF (Argon2id parameters), a wrapping cipher, a versioning byte, and a textual armor format --- each a non-trivial decision that ages on a different clock than the wire format itself. The lesson from \texttt{age}'s ecosystem is that this work belongs to plugins (\texttt{age-plugin-yubikey}, \texttt{age-plugin-tpm}, \texttt{age-plugin-se}, \texttt{age-plugin-fido2}). PQCrypt ships the plugin protocol so plugin authors can build immediately, plus a reference plugin (\texttt{pqcrypt-plugin-file}, \cref{sec:plugin-protocol:reference}) so the v1 spec is self-contained.

\subsection{Wire format}

The plugin protocol uses stdin/stdout binary frames modeled on the C2SP \texttt{age-plugin} specification \citep{c2sp-age-plugin}. The PQCrypt library invokes \texttt{pqcrypt-plugin-\textless{}name\textgreater} via \texttt{exec}, passes envelope-level data through framed messages, and receives wrapped/unwrapped key material in return.

Frames are length-prefixed: 4-byte big-endian length, then the frame body. Frame types are identified by a leading single-byte tag. v1 frame types:

\begin{itemize}
\item \texttt{RECIPIENT\_STANZA} (\code{0x01}): declares the plugin's recipient identifier in URI form \texttt{pqck1\textless{}plugin-name\textgreater{}1\textless{}base64url-payload\textgreater}.
\item \texttt{WRAP\_KEY} (\code{0x02}): library $\to$ plugin, requesting wrap of a per-envelope file-key.
\item \texttt{UNWRAP\_KEY} (\code{0x03}): library $\to$ plugin, requesting unwrap of a previously-wrapped file-key.
\item \texttt{OK} (\code{0x10}): plugin $\to$ library, success with body payload.
\item \texttt{ERROR} (\code{0x11}): plugin $\to$ library, failure with optional explanatory string.
\end{itemize}

The full state machine, including handshake and pipelining rules, is defined in \texttt{pqcrypt/plugin/pqcrypt-plugin-protocol.md}.

\subsection{Reference plugin: \texttt{pqcrypt-plugin-file}}
\label{sec:plugin-protocol:reference}

\texttt{pqcrypt-plugin-file} stores a passphrase-wrapped ML-KEM-1024 private key on disk. Construction:

\begin{enumerate}
\item Read the passphrase from the terminal (no echo) or from \code{\$PQCRYPT\_PASSPHRASE}.
\item Derive a 32-byte symmetric key: \code{K\_wrap = Argon2id(passphrase, salt, m=64MiB, t=3, p=1, 32 bytes)} \citep{rfc9106}.
\item Wrap the ML-KEM-1024 private key ($\sim$3168 B) under AES-256-GCM: \code{wrapped = AES-256-GCM-Encrypt(K\_wrap, nonce, sk)} where \code{nonce} is 12 fresh CSPRNG bytes.
\item Serialize: \texttt{pqck-key-v1\textbackslash{}n\textless{}base64url(salt || nonce || wrapped || tag)\textgreater\textbackslash{}n}.
\end{enumerate}

The reference plugin is not part of the v1 normative wire format; only the plugin \emph{protocol} is. Other plugins (\texttt{pqcrypt-plugin-yubikey}, hypothetical) may wrap keys however they want.

\subsection{What the protocol does \emph{not} cover}

\begin{itemize}
\item Key custody (HSM, secure enclave, TPM): a job for plugin authors.
\item Key rotation: a higher-level concern; PQCrypt envelopes are self-contained and do not encode rotation lineage.
\item Plugin trust: PQCrypt invokes whichever plugin binary the caller's PATH resolves; the caller is responsible for trusting that binary.
\end{itemize}
